% Artículo VII — primera redacción editorial, 10 de septiembre de 2026.
% Apertura integrada con las demostraciones de la revisión 06.
% Procedencia y condiciones: PLAN_RESIDENCIAS_DEMOSTRATIVAS.md.
% Las referencias internas remiten a demostraciones incluidas en el cuerpo.

\begin{center}\bfseries Resumen\end{center}
\begingroup
\fontsize{10}{12}\selectfont
\setlength{\parskip}{2pt plus 1pt minus .5pt}
\noindent
Se determina la amplitud torsional como funcional explícito del estado
material mediante la acción espinorial, la eliminación de la contorsión
y la variación métrica. Se construye un sector material en el que esa
amplitud es positiva y se conserva. Su contribución energética produce
una dilución \(a^{-6}\) y, en la realización homogénea plana con
\(\Lambda_{\mathrm{eff}}=0\), densidades positivas y \(w<1\), un umbral
de rebote único y transversal. Para polvo se obtiene una solución exacta
para todo tiempo, con curvatura finita y geodésicas causales de
Levi--Civita completas en ambas direcciones.

La construcción parte de un núcleo formal denominado Holografía Modular
Triádica, en adelante HMT. \emph{All Possible Paths}, en adelante APP,
reúne caminos aritméticos sobre \((\mathbb Z/9\mathbb Z)^2\), con
adyacencia dada por el grafo de Cayley aditivo de generadores cardinales
y realización celular toroidal, distinguiendo las operaciones y las
evaluaciones de suma y producto;
TRIT es la unidad de firma orientada \(\{-1,0,+1\}\) que determina los
regímenes locales; el \emph{Topological Phase Kernel}, en adelante TPK,
compone la dinámica de transporte y memoria. El estado enriquecido
conserva hojas, orientación, residuo, acarreo e historia durante el
retorno de fase. La Mecánica Dimensional expresa sus resultados
internos mediante magnitudes y observables físicos, preservando la
procedencia estructural de los acoplamientos.

La normalización gravitatoria se desarrolla previamente mediante la
composición longitudinal \(L=(5759/23040)\alpha^{16}L_*\) y su
métrica radial positiva. En la realización definida por esas operaciones
y por la identificación del radio gravitatorio reducido, la reciprocidad
con la sección de acción fija \(G=c^3L^2/\hbar_{\rm ret}\) en todos
los modos admisibles. La demostración conserva el archivo completo y
establece las condiciones de reducción estática y dinámica de la memoria.
La elección dimensional \(L_*=1\,\mathrm m\), las reglas longitudinal
y métrica y la identificación radial se declaran en el propio enunciado.

El transporte de rutas admite una realización de Cartan compatible con
composición, inversión y refinamiento. Sobre ella se distinguen dos
acciones materiales: la extensión minimizante de la pantalla de memoria
y la realización espinorial. Cada una conserva su propia variación
respecto de la conexión. En el sector espinorial, afín en la contorsión,
la inversión de los operadores de Holst y torsión determina la corrección
efectiva; para la extensión mínima se conserva la dependencia no lineal
de la corriente y se demuestra un criterio local mediante el Hessiano.
La corriente, el observable cuártico y su normalización explican el
valor de la amplitud, además de su papel en la evolución cosmológica.

La persistencia del rebote bajo perturbaciones se demuestra mediante
cotas locales diferenciadas, mientras la completitud causal corresponde
a la solución global explícita de polvo. Los balances residuales, la
información de horizonte y las operaciones del observador conservan las
condiciones de fuentes, acoplamientos y cartas que utilizan. La respuesta
relacional del estado conjunto se enlaza con la propagación avanzada y
retardada mediante operadores de Green para fuentes de soporte finito,
asociados a la diferencia covariante
construida a partir del transporte de rutas y una condición de
compatibilidad en la frontera. El resultado relaciona la memoria
generada, su realización geométrica y las consecuencias dinámicas e
informacionales de esa realización.

La composición entre acción, incremento areal de Barbero--Immirzi y
horizonte determina la conversión
\(T_{\rm cos}(L_\alpha)/\Theta_{\rm clk}=\log3/(2\pi)\).
El Hamiltoniano de respuesta angular proporciona un balance finito de
energía libre mediante la entropía relativa y una primera ley a radio
fijo. Una sección termométrica relativa caracteriza, además, el
transductor positivo y el criterio para identificar otro lector sobre
esa misma sección.
\par\endgroup
\clearpage

\section{Introducción}
\label{vii:sec:introduccion}

La cuestión gravitatoria que organiza este trabajo es cómo una estructura
que conserva la historia de sus operaciones puede determinar una
corriente material y, mediante ella, una contribución a la evolución del
espacio-tiempo. La respuesta se construye desde el transporte de rutas
hasta la variación de la acción: el resultado central es una amplitud
torsional explícita, acompañada de un sector de positividad conservada y
de una solución cosmológica global en la realización de polvo. La tesis
vincula la memoria aritmética con una respuesta geométrica determinada por
operaciones; el retorno de fase conserva una periodicidad observable
mientras el estado acumula orientación, acarreo y profundidad.

Se introduce para ello un núcleo formal matemático denominado Holografía
Modular Triádica, en adelante HMT. \emph{All Possible Paths}, en adelante
APP, es una estructura aritmética de caminos sobre
\((\mathbb Z/9\mathbb Z)^2\). Su adyacencia se describe mediante el grafo
de Cayley aditivo con generadores cardinales \(\pm(1,0)\) y
\(\pm(0,1)\), y su realización celular periódica es toroidal. La acumulación
aditiva, la composición multiplicativa y sus evaluaciones mantienen
funciones distintas dentro de esa estructura. TRIT es la unidad de firma
orientada \(\{-1,0,+1\}\), que tipa el régimen local de la operación.
El \emph{Topological Phase Kernel}, en adelante TPK, constituye su
dinámica de selección, transporte, memoria y prolongación. El estado
enriquecido y la estructura discreta conjunta del continuo conservan las
hojas y los datos residuales a través de esas operaciones. La Mecánica
Dimensional realiza sus salidas en secciones de acción, magnitudes
geométricas y observables físicos; los valores internos transmiten a
cada etapa las relaciones que los determinan.

La cuestión central es cómo la memoria transportada llega a intervenir
en la evolución gravitatoria. La dilución, el rebote y su persistencia
describen las consecuencias cosmológicas de una amplitud que antes debe
construirse desde la acción material. La primera relación se expresa
mediante dos lectores de una misma ruta,
\begin{equation}
 \frac{a_n}{a_{\mathrm{ref}}}=\varphi^{n/3},
 \qquad M_n=\varphi^{-2n},
 \qquad
 M_n=\left(\frac{a_n}{a_{\mathrm{ref}}}\right)^{-6}.
 \label{vii:eq:programa-dilucion}
\end{equation}
Aquí \(a_{\mathrm{ref}}>0\) fija la normalización del factor de escala.
La potencia expresa la relación entre el crecimiento de la realización
espacial y el carácter cuadrático del lector de memoria. La contribución
gravitatoria requiere, además, determinar la corriente y la amplitud que
acompañan a esa dependencia de escala. Ambas construcciones ocupan una
sección propia antes de introducirlas en las ecuaciones cosmológicas.

El primer movimiento del artículo reúne construcción y geometría. La
descomposición del transporte distingue memoria simétrica y circulación
orientada. La representación del grupoide de rutas en geometría de Cartan
se desarrolla con composición, inversión, covariancia y refinamiento. La
acción, el parámetro de Barbero--Immirzi, la constante gravitatoria y la
unidad areal conservan sus antecedentes estructurales; su reutilización
transmite las condiciones que individualizan cada sección.

En particular, la sección~\ref{g26:sec:rigidez-radial-es} construye
primero la longitud \(L\), compone su lector positivo con el mismo
carácter que determina la energía y obtiene el coeficiente gravitatorio.
El reloj circular recibe entonces
\(t_0=\pi_{\rm HMT}L/(54c)\); su radio \(R_{108}\) coincide con
\(L\) en esta realización. El orden de la demostración distingue la
composición longitudinal, la respuesta operatoria y la evaluación en
unidades físicas. Las reglas longitudinal y métrica son composiciones
explicitadas en este desarrollo; el teorema prueba la rigidez de las
realizaciones que las conservan, con la identificación radial declarada.

El segundo movimiento desarrolla amplitud, evolución y persistencia.
La corriente de pantalla se lleva a un diferencial respecto de la
conexión mediante la incidencia transportada y la extensión minimizante.
La realización espinorial se desarrolla, a su vez, desde las matrices
internas de Clifford y su acción material. La acción de extensión
minimizante y la acción espinorial son dos funcionales distintos: la
primera depende en general no linealmente de la conexión, mientras la
segunda es afín en la contorsión cuando se fijan cotetrada y campo
espinorial. Sus variaciones se realizan sobre los argumentos propios de
cada funcional; su posible composición exige conservar esa dependencia.
La normalización tracial
fija la lectura de densidad correspondiente y se mantiene diferenciada de
una suma extensiva. La variación de la conexión determina la respuesta
torsional de la acción mínima. Al evolucionar la cotetrada y los campos
materiales, cambia también la torsión algebraicamente asociada a su
corriente de espín.

La amplitud del sector espinorial queda determinada por
\begin{equation}
 s_0[\varrho,V_c]
 =\frac{3\kappa c^2\hbar^2}{16}
   \frac{\gamma^2}{1+\gamma^2}
   \frac{\operatorname{Tr}
    (\varrho\widehat{\mathscr Q}_{\rm sp})}{V_c^2},
 \qquad \kappa=\frac{8\pi G}{c^4}.
 \label{vii:eq:programa-amplitud}
\end{equation}
Aquí \(\varrho\) es el estado material normalizado, \(V_c\) el volumen
comóvil y \(\widehat{\mathscr Q}_{\rm sp}\) el observable cuártico
construido a partir de la corriente. El
teorema~\ref{vii:thm:amplitud-material-explicita} deriva la fórmula y
el corolario~\ref{vii:cor:dominio-positivo-conservado} exhibe estados
en los que la amplitud es positiva y permanece constante. La dependencia
del segundo momento conserva información del estado incluso cuando la
media de la corriente se anula. El dato material y su volumen fijan una
realización; las constantes de acoplamiento proceden de las secciones
estructurales construidas previamente.

La eliminación de la torsión se trata también para la acción de
extensión mínima, cuya dependencia respecto de la conexión es, en
general, no lineal. Su diferencial determina la ecuación estacionaria;
el Hessiano controla la continuación local y una identidad integral
determina la acción efectiva. Esta formulación permite conservar la
dependencia completa del mínimo al efectuar las variaciones.

La formulación del rebote conserva las densidades positivas, el signo de
la contribución torsional y la desigualdad de exponentes de la realización
homogénea. Su persistencia local se demuestra mediante dos controles
separados: una cota uniforme de la perturbación conserva el cambio de
signo en los extremos de un intervalo; una cota uniforme de su derivada
conserva la monotonía y la transversalidad. Las realizaciones de los
selectores cosmológicos se enuncian con sus propios dominios, junto a los
invariantes y balances que producen.

La solución plana de polvo con \(\Lambda_{\mathrm{eff}}=0\)
proporciona además un desarrollo global:
\[
 a(t)=a_{\rm b}\left(1+
       \frac{(t-t_{\rm b})^2}{\tau^2}\right)^{1/3},
 \qquad
 a_{\rm b}^3=\frac{s_0}{\rho_0},\qquad
 \tau^2=\frac{s_0}{6\pi G\rho_0^2},
 \quad c=1.
\]
La prueba establece regularidad de la curvatura y completitud de las
geodésicas causales de Levi--Civita en esa realización homogénea.
Su integración y la persistencia local
del umbral describen dos propiedades complementarias, con sus
respectivos dominios explícitos.

La conservación del tensor residual se formula en la carta métrica
seleccionada, con acoplamientos constantes sobre ella y una fuente
material covariantemente conservada. Su descomposición en traza y parte
sin traza permite identificar el intercambio entre ambos sectores;
la traza constante caracteriza el caso de conservación separada.

El tercer movimiento estudia horizontes y respuesta relacional. El estado
reducido, los registros accesibles y la purificación permiten describir
qué información representa una lectura y qué información permanece en las
correlaciones del estado conjunto. La definición operacional del
observador comprende orientación y trivialización del transporte,
restricción a observables y proyección al subsistema. La respuesta
autoinercial se conecta con esas operaciones mediante la factorización
por la traza global y el defecto geométrico de proyección.

\Needspace{24\baselineskip}
Los balances de horizonte conservan igualmente las convenciones de
área, energía, temperatura y estado que utilizan.
Las secciones~\ref{viii:subsec:conversion-termica-calendario-horizonte}--\ref{viii:subsec:seccion-termometrica-registro}
componen el cuanto areal con la energía de decisión del calendario y
obtienen la conversión entre las temperaturas cronológica y cosmológica.
El estado reducido conserva las ocupaciones de los dos sectores
angulares. A radio fijo, su respuesta energética tiene Hamiltoniano
\(\mathsf H_R=\tau_R\Delta_0\mathsf D_\partial\); sus variaciones
producen la primera ley, mientras las diferencias finitas retienen
el término de entropía relativa. Este desarrollo construye una sección
termométrica relativa y establece en ella una prueba de igualdad entre
mapas lineales de conversión. La procedencia sectorial permanece
accesible tanto en el área como en el coste energético.
La propagación avanzada y retardada para fuentes de soporte finito se
compone sobre el transporte unitario de rutas. La igualdad de sus lecturas en la frontera se
expresa mediante una restricción lineal de fuentes; su proyección
determina la respuesta variacional en ese dominio.

La unidad del argumento reside en la relación entre fuente, corriente,
amplitud y geometría. La memoria conserva información que la lectura de
fase no agota; su realización material produce un segundo momento capaz
de contribuir a la dinámica incluso con corriente media nula. La
evolución cosmológica expresa esa contribución a escala global en el
sector resuelto, y las operaciones del observador precisan qué parte de
la información permanece accesible. La estructura, la ecuación y su
solución quedan así enlazadas dentro de realizaciones con dominios
explícitos.
